Neighbour selection of the two adaptive rules on one node with seven candidate neighbours A--G, raw scores $s = (2.1, 1.9, 1.8, 0.3, 0.1, -0.2, -0.5)$ in (a) and the same scores divided by 10 in (b). For each candidate, the left bar is coloured if the Jenks rule keeps it and the right bar if entmax-1.5 keeps it (grey: dropped). Jenks splits the sorted scores into two natural-breaks classes (dashed line) and keeps the upper class; entmax-1.5 assigns $p_j = [s_j/2 - \tau]_+^2$ with $\sum_j p_j = 1$ and keeps every candidate with $p_j > 0$, i.e.\ $s_j > 2\tau$ (dotted line); the numbers above the bars are $p_j$. The Jenks split depends only on the shape of the score row, so it keeps three neighbours in both panels, whereas entmax-1.5 keeps three in (a) and all seven once the scores are flatter (b). The kept sets were computed by the model's own adjacency module on these scores (no training involved).
